\documentclass[a4paper,fontsize=14pt,fleqn]{scrartcl}%fleqn pour aligner les equations à gauche

\usepackage{../../Styles/Eval}

\newcommand{\chnum}{Chapitre 4}
\newcommand{\chtitle}{Avancement et oxydoréduction}
\newcommand{\dtype}{DS3}

\begin{document}
\maketitle

\section{Nitroglycérine}

La nitroglycérine est utilisée dans la fabrication de la dynamite. Le processus de sa fabrication date de 1860. Beaucoup de chimistes s'y sont essayés, certains y ont laissé leur vie comme le frère d'Alfred Nobel, son inventeur. La nitroglycérine est en effet extrêmement instable et peut exploser au moindre choc ou à une élévation de température. Elle peut être fabriquée en laboratoire dans des condictions très particulières en faisant agit \SI{1,00}{L} de glycérol \ce{C3H5(OH)3} avec \SI{1,00}{L} d'acide nitrique (\ce{H+}, \ce{NO3-}) selon l'équation : $$ \ce{C3H5(OH)3(l) + 3NO3(l) -> C3H5(NO3)3(l) + 3H2O(l)}$$

\begin{enumerate}
    \item Calculer les quantités initiales des réactifs.
    \item Donner la définition du réactif limitant.
    \item En vous aidant d'un tableau d'avancement, déterminer la nature du réactif limitant.
    \item La réaction est-t-elle totale ?
    \item Calculer la masse de nitroglycérine ainsi fabriquée.
    \item Calculer le volume de l'eau formée lors de la réaction.
\end{enumerate}

\begin{data}
    \begin{itemize}
        \item Glycérol : $\rho = \SI{1,26}{\gram\per\centi\meter\cubed}$ ; $M = \SI{92,0}{\gram\per\mole}$
        \item Acide nitrique : $\rho = \SI{1,51}{\gram\per\centi\meter\cubed}$ ; $M = \SI{63,0}{\gram\per\mole}$
        \item Nitroglycérine : $\rho = \SI{1,60}{\gram\per\centi\meter\cubed}$ ; $M = \SI{227,0}{\gram\per\mole}$
        \item Eau : $\rho = \SI{1,00}{\gram\per\centi\meter\cubed}$ ; $M = \SI{18,0}{\gram\per\mole}$
    \end{itemize}
\end{data}

\newpage
\section{Piles électrochimiques}
Georges Leclanché a été un des précurseurs des piles en proposant en 1868 une pile saline zinc/diodyde de manganèse. De nombreux autres types de piles électrochimiques ont depuis été inventées, comme la pile alcaline. La différence entre la pile saline et la pile alcaline réside dans l'électrolyte (corps entre les électrodes). Ainsi, il peut exister des piles zinc/diodyde de manganèse saline ou alcalines. Celles-ci mettent en jeu les couples \ce{ZnO(s)}/\ce{Zn(s)} et \ce{MnO2(s)}/\ce{Mn2O3(s)} en milieu basique.

\begin{enumerate}
    \item Déterminer les demi équations des couples suivants : \begin{enumerate}
        \item \ce{Zn^2+(aq)}/\ce{Zn(s)}
        \item \ce{MnO4-(aq)}/\ce{MnO2(s)}
        \item \ce{SO4^{2-}(aq)}/\ce{SO3^{2-}(aq)}
    \end{enumerate}
    \item Écrire les demi-équations des couples \ce{ZnO(s)}/\ce{Zn(s)} et \ce{MnO2(s)}/\ce{Mn2O3(s)}.
    \item En déduire l'équation de la réaction.
\end{enumerate}

On réalise une pile saline zinc/dioxyde de manganèse en laboratoire. Pour celà on dispose de \SI{250}{g} de zinc solide et de \SI{800}{g} de dioxyde de manganèse solide.
\begin{enumerate}[resume]
    \item Etablir un tableau d'avancement de la réaction.
    \item Déterminer la nature du réactif limitant.
    \item Calculer la masse de produit formé.
    \item La masse de la pile change t-elle au cours de la réaction ? Justifier.
\end{enumerate}

\begin{data}
    \begin{itemize}
        \item Masse molaire du zinc solide : $M(\ce{Zn}) = \SI{65,4}{\gram\per\mole}$
        \item Masse molaire du dioxyde de manganèse solide : $M(\ce{MnO2}) = \SI{86,9}{\gram\per\mole}$
        \item Masse molaire de l'oxyde de zinc solide : $M(\ce{ZnO}) = \SI{81,4}{\gram\per\mole}$
        \item Masse molaire du trioxyde de manganèse solide : $M(\ce{Mn2O3}) = \SI{157,8}{\gram\per\mole}$
    \end{itemize}
\end{data}
\end{document}

